\documentclass[../../../uem_tok_q.tex]{subfiles}

\begin{document}
    座標平面上の点$\mathrm{A}(0,\;0)$，$\mathrm{B}(0,\;1)$，%
    $\mathrm{C}(1,\;1)$，$\mathrm{D}(1,\;0)$を考える．%
    実数 $0<t<1$ に対して，線分$\mathrm{A}\mathrm{B}$，%
    $\mathrm{B}\mathrm{C}$，$\mathrm{C}\mathrm{D}$を $t:(1-t)$ に内分する点をそれぞれ%
    $\mathrm{P}_t$，$\mathrm{Q}_t$，$\mathrm{R}_t$とし，%
    線分$\mathrm{P}_t\mathrm{Q}_t$，$\mathrm{Q}_t\mathrm{R}_t$を %
    $t:(1-t)$ に内分する点をそれぞれ$\mathrm{S}_t$，$\mathrm{T}_t$とする．%
    さらに，線分$\mathrm{S}_t\mathrm{T}_t$を $t:(1-t)$ に内分する点を$\mathrm{U}_t$とする．%
    また，点Aを$\mathrm{U}_0$，点Dを$\mathrm{U}_1$とする．

    \begin{enumerate}
        \item 点$\mathrm{U}_t$の座標を求めよ．
        \item $t$が $0\leqq t\leqq1$ の範囲を動くときに点$\mathrm{U}_t$が描く曲線と，%
            線分ADで囲まれた部分の面積を求めよ．
        \item $a$を $0<a<1$ を満たす実数とする．$t$が $0\leqq t\leqq a$ の範囲を動くときに%
            点$\mathrm{U}_t$が描く曲線の長さを，$a$の多項式の形で求めよ．
    \end{enumerate}
\end{document}
